Include a header paragraph specific to your product here. Safety requirements address hazards to people, property, and the environment arising from developing, operating, or maintaining the product: exposure to toxic materials, sharp edges, breakable glass, laser/IR eye exposure, electrical shock, stored mechanical or chemical energy, moving machinery, etc. The examples below are common to the CSE Senior Design laboratory; retain those that apply and add product-specific safety requirements.

\subsection{SA-01: Laboratory equipment lockout/tagout (LOTO) procedures}
\subsubsection{Description}
Any fabrication equipment used in the development of the project shall be used in accordance with OSHA standard LOTO procedures. Locks and tags are installed on all equipment items that present use hazards, and only the course instructor or designated teaching assistants may remove a lock. All locks will be immediately replaced once the equipment is no longer in use.
\subsubsection{Rationale}
Prevents injury from unexpected energization or start-up of laboratory equipment during service or setup.
\subsubsection{Source}
CSE Senior Design laboratory policy.
\subsubsection{Priority}
Critical
\subsubsection{Verification Method}
Inspection (locks/tags present and controlled) and Demonstration (procedure followed).
\subsubsection{Constraints}
Equipment usage, due to lock removal policies, is limited to the availability of the course instructor and designated teaching assistants.
\subsubsection{Applicable Standards}
OSHA 29 CFR 1910.147 --- The control of hazardous energy (lockout/tagout).

\subsection{SA-02: National Electric Code (NEC) wiring compliance}
\subsubsection{Description}
Any electrical wiring shall be completed in compliance with all requirements specified in the National Electric Code, including wire runs, insulation, grounding, enclosures, over-current protection, and all other specifications.
\subsubsection{Rationale}
Reduces risk of fire and electric shock in fabricated electrical assemblies.
\subsubsection{Source}
CSE Senior Design laboratory policy.
\subsubsection{Priority}
Critical
\subsubsection{Verification Method}
Inspection of wiring against NEC provisions.
\subsubsection{Constraints}
High-voltage power sources, as defined in NFPA 70, will be avoided as much as possible to minimize hazards.
\subsubsection{Applicable Standards}
NFPA 70 (National Electrical Code).

\subsection{SA-03: Robotic manipulator safety}
\subsubsection{Description}
Robotic manipulators, if used, shall either be housed in a compliant lockout cell with all required safety interlocks, or be certified as a ``collaborative'' unit from the manufacturer.
\subsubsection{Rationale}
Protects personnel from injury by industrial robot motion.
\subsubsection{Source}
CSE Senior Design laboratory policy.
\subsubsection{Priority}
Critical
\subsubsection{Verification Method}
Inspection of interlocks/certification; Demonstration of safe operation.
\subsubsection{Constraints}
Collaborative manipulators are preferred over non-collaborative units to minimize hazards. Sourcing and use of any required safety interlock mechanisms is the responsibility of the engineering team.
\subsubsection{Applicable Standards}
ANSI/RIA R15.06 --- Industrial Robots and Robot Systems; RIA TR15.606 --- Collaborative Robots.
